% GENERATED FILE. Edit canonical lessons, then run npm run generate:books.
% canonical-source-hash: fnv1a64:cb5d268780352f57
% canonical-lessons: KA-C168-hane, KA-C168-gantalu, KA-C168-himmadi, KA-C168-tode, KA-C168-angai

\chapter{Forehead, Throat, Heel, Thigh, Palm}
\label{ch:ka-forehead-throat-heel-thigh-palm}
\hlchaptermodality{168}

\begin{chapteropening}
\textbf{By the end of this chapter:} \emph{I can say the forehead, the throat, the heel, the thigh and the palm.}

Complete the last lesson of chapter 168: I can say the forehead, the throat, the heel, the thigh and the palm.
\end{chapteropening}

\section[\texorpdfstring{\kn{ಹಣೆ} (haṇe)}{ಹಣೆ (haṇe)}]{\kn{ಹಣೆ} (haṇe) — the forehead}

\label{lesson:KA-C168-hane}



\begin{quote}
Before the new one: say the Kannada for to burst, then the Kannada for to snore.
\end{quote}

\subsection*{You'll want to know: \kn{ಹಣೆ}}
\textbf{\kn{ಹಣೆ}} — \emph{haṇe} — \textquotedblleft{}the forehead\textquotedblright{}.

\begin{grammarlens}[title={What you've built}]
One more word for this chapter.
\end{grammarlens}

\subsection*{Your turn}
\begin{itemize}
\raggedright
  \item \emph{Say it:} \emph{haṇe}
  \item \emph{Say it:} \emph{haṇe}, once more
  \item \emph{Say it:} say \emph{haṇe}
  \item \emph{Recall:} say the Kannada for to burst, then the Kannada for to snore, then say \emph{haṇe} again
\end{itemize}

\begin{tcolorbox}[breakable,colback=teal!4,colframe=teal!35!black,title={Before you move on}]
What does \kn{ಹಣೆ} mean? (\textquotedblleft{}The forehead\textquotedblright{}.) Say it once more.
\end{tcolorbox}
\section[\texorpdfstring{\kn{ಗಂಟಲು} (gaṇṭalu)}{ಗಂಟಲು (gaṇṭalu)}]{\kn{ಗಂಟಲು} (gaṇṭalu) — the throat}

\label{lesson:KA-C168-gantalu}



\begin{quote}
Before the new one: say the Kannada for to snore, then the Kannada for the forehead.
\end{quote}

\subsection*{You'll want to know: \kn{ಗಂಟಲು}}
\textbf{\kn{ಗಂಟಲು}} — \emph{gaṇṭalu} — \textquotedblleft{}the throat\textquotedblright{}.

\begin{grammarlens}[title={What you've built}]
One more word for this chapter.
\end{grammarlens}

\subsection*{Your turn}
\begin{itemize}
\raggedright
  \item \emph{Say it:} \emph{gaṇṭalu}
  \item \emph{Say it:} \emph{gaṇṭalu}, once more
  \item \emph{Say it:} say \emph{gaṇṭalu}
  \item \emph{Recall:} say the Kannada for to snore, then the Kannada for the forehead, then say \emph{gaṇṭalu} again
\end{itemize}

\begin{tcolorbox}[breakable,colback=teal!4,colframe=teal!35!black,title={Before you move on}]
What does \kn{ಗಂಟಲು} mean? (\textquotedblleft{}The throat\textquotedblright{}.) Say it once more.
\end{tcolorbox}
\section[\texorpdfstring{\kn{ಹಿಮ್ಮಡಿ} (himmaḍi)}{ಹಿಮ್ಮಡಿ (himmaḍi)}]{\kn{ಹಿಮ್ಮಡಿ} (himmaḍi) — the heel}

\label{lesson:KA-C168-himmadi}



\begin{quote}
Before the new one: say the Kannada for the forehead, then the Kannada for the throat.
\end{quote}

\subsection*{You'll want to know: \kn{ಹಿಮ್ಮಡಿ}}
\textbf{\kn{ಹಿಮ್ಮಡಿ}} — \emph{himmaḍi} — \textquotedblleft{}the heel\textquotedblright{}.

\begin{grammarlens}[title={What you've built}]
One more word for this chapter.
\end{grammarlens}

\subsection*{Your turn}
\begin{itemize}
\raggedright
  \item \emph{Say it:} \emph{himmaḍi}
  \item \emph{Say it:} \emph{himmaḍi}, once more
  \item \emph{Say it:} say \emph{himmaḍi}
  \item \emph{Recall:} say the Kannada for the forehead, then the Kannada for the throat, then say \emph{himmaḍi} again
\end{itemize}

\begin{tcolorbox}[breakable,colback=teal!4,colframe=teal!35!black,title={Before you move on}]
What does \kn{ಹಿಮ್ಮಡಿ} mean? (\textquotedblleft{}The heel\textquotedblright{}.) Say it once more.
\end{tcolorbox}
\section[\texorpdfstring{\kn{ತೊಡೆ} (toḍe)}{ತೊಡೆ (toḍe)}]{\kn{ತೊಡೆ} (toḍe) — the thigh}

\label{lesson:KA-C168-tode}



\begin{quote}
Before the new one: say the Kannada for the throat, then the Kannada for the heel.
\end{quote}

\subsection*{You'll want to know: \kn{ತೊಡೆ}}
\textbf{\kn{ತೊಡೆ}} — \emph{toḍe} — \textquotedblleft{}the thigh\textquotedblright{}.

\begin{grammarlens}[title={What you've built}]
One more word for this chapter.
\end{grammarlens}

\subsection*{Your turn}
\begin{itemize}
\raggedright
  \item \emph{Say it:} \emph{toḍe}
  \item \emph{Say it:} \emph{toḍe}, once more
  \item \emph{Say it:} say \emph{toḍe}
  \item \emph{Recall:} say the Kannada for the throat, then the Kannada for the heel, then say \emph{toḍe} again
\end{itemize}

\begin{tcolorbox}[breakable,colback=teal!4,colframe=teal!35!black,title={Before you move on}]
What does \kn{ತೊಡೆ} mean? (\textquotedblleft{}The thigh\textquotedblright{}.) Say it once more.
\end{tcolorbox}
\section[\texorpdfstring{\kn{ಅಂಗೈ} (aṅgai)}{ಅಂಗೈ (aṅgai)}]{\kn{ಅಂಗೈ} (aṅgai) — the palm (of the hand)}

\label{lesson:KA-C168-angai}



\begin{quote}
Before the new one: say the Kannada for the heel, then the Kannada for the thigh.
\end{quote}

\subsection*{You'll want to know: \kn{ಅಂಗೈ}}
\textbf{\kn{ಅಂಗೈ}} — \emph{aṅgai} — \textquotedblleft{}the palm (of the hand)\textquotedblright{}.

\begin{grammarlens}[title={What you've built}]
One more word for this chapter.
\end{grammarlens}

\subsection*{Your turn}
\begin{itemize}
\raggedright
  \item \emph{Say it:} \emph{aṅgai}
  \item \emph{Say it:} \emph{aṅgai}, once more
  \item \emph{Say it:} say \emph{aṅgai}
  \item \emph{Recall:} say the Kannada for the heel, then the Kannada for the thigh, then say \emph{aṅgai} again
\end{itemize}

\begin{tcolorbox}[breakable,colback=teal!4,colframe=teal!35!black,title={Before you move on}]
What does \kn{ಅಂಗೈ} mean? (\textquotedblleft{}The palm (of the hand)\textquotedblright{}.) Say it once more.
\end{tcolorbox}
